\section{Learned Bounded Refinement}

The method keeps Tensor-OMP support selection and least-squares reconstruction explicit. Learning controls only how far the estimate moves from a selected grid atom toward a deterministic local-search candidate. Figure~\ref{fig:architecture} separates the two learned implementations from the repeated-refinement depth diagnostic; they share a bounded interpolation principle but are not one integrated network.

\begin{figure*}[t]
  \centering
  \includegraphics[width=\textwidth]{v2_3r_refinement_paths.pdf}
  \caption{Implemented refinement paths and evidence scope. The local three-dimensional experiment learns one scalar interpolation parameter. The CDL experiment uses a separate 10--16--2 feature controller. A4 repeats deterministic local searches and applies a scalar plateau gate; it is a depth diagnostic, not the depth of the CDL controller.}
  \label{fig:architecture}
\end{figure*}

\subsection{Tensor-OMP Support and Bounded Interpolation}

Let $\mathbf{b}_{0,\ell}$ denote the coarse tensor-grid coordinates selected for component $\ell$, and let $\mathbf{b}_{\star,\ell}$ be the candidate returned by deterministic bounded local search around that support. The learned coordinates are
\begin{equation}
  \widehat{\mathbf{b}}_{\ell}
  =\mathbf{b}_{0,\ell}
  +\boldsymbol{\alpha}\odot
   (\mathbf{b}_{\star,\ell}-\mathbf{b}_{0,\ell}),
  \qquad 0<\alpha_d<1.2,
  \label{eq:bounded_interpolation}
\end{equation}
where a scalar $\alpha$ is broadcast over axes in the local implementation and axis-specific scales are used in the CDL implementation. Given the design matrix $\mathbf{A}(\widehat{\mathbf{B}})$ formed from all refined atoms, coefficients and reconstruction are recomputed as
\begin{equation}
  \widehat{\boldsymbol{\beta}}
  =\arg\min_{\boldsymbol{\beta}}
   \|\mathbf{y}-\mathbf{A}(\widehat{\mathbf{B}})\boldsymbol{\beta}\|_2^2,
  \qquad
  \widehat{\mathbf{y}}=\mathbf{A}(\widehat{\mathbf{B}})\widehat{\boldsymbol{\beta}}.
  \label{eq:ls_reconstruction}
\end{equation}
Thus the learned output remains an inspectable physical support followed by an LS estimate.

\subsection{Local Three-Dimensional Scalar Model}

For the $128\times16\times32$ angle--delay--Doppler experiment, the trainable model contains one unconstrained scalar $a$. Its bounded interpolation scale is
\begin{equation}
  \alpha=1.2\,\sigma(a),
  \label{eq:scalar_alpha}
\end{equation}
and $a$ is optimized through the differentiable LS reconstruction to minimize linear-scale measurement NMSE. The coarse support and local candidate are prepared before this optimization. Consequently, the G2 result tests whether a learned global step length improves a deterministic local candidate; it does not test a deep stack of trainable Tensor-OMP layers.

\subsection{Feature-Conditioned CDL Controller}

The CDL path uses a separate controller because a single shared delay--angle scale produced strong axis coupling. A standardized feature vector $\mathbf{z}\in\mathbb{R}^{10}$ records support density, selected spectral energy, grid residual energy, delay/angle displacement statistics, support uniqueness, and design-matrix conditioning. The controller is
\begin{equation}
  \begin{bmatrix}\alpha_\tau\\ \alpha_\theta\end{bmatrix}
  =1.2\,\sigma\!\left(
    \mathbf{W}_2\tanh(\mathbf{W}_1\mathbf{z}+\mathbf{b}_1)+\mathbf{b}_2
  \right),
  \label{eq:feature_controller}
\end{equation}
with dimensions $10\!\rightarrow\!16\!\rightarrow\!2$, for 210 trainable parameters. At the four-element ULA Nyquist bin, the angle coordinate is locked to its coarse value to avoid changing a valid endfire representation through an aliased continuous update.

For each training sample, the current predicted set is matched to the reference set with Hungarian assignment. The optimized objective is
\begin{equation}
  \mathcal{L}
  =\operatorname{NMSE}(\widehat{\mathbf{y}},\mathbf{y}_{\rm clean})
  +\lambda_p\bigl[
      \operatorname{SL1}_{0.5}(\widehat{\boldsymbol{\tau}},\boldsymbol{\tau})
      +\operatorname{SL1}_{0.25}(\mathbf{e}_\theta/30)
    \bigr],
  \label{eq:training_loss}
\end{equation}
where $\mathbf{e}_\theta$ is the minimum valid projected-angle error over the ULA alias candidates. Assignment and the alias-aware angle calculation are recomputed from the current prediction. This makes supervision permutation-invariant while keeping the physical objective aligned with test-time metrics.

\subsection{Separate Repeated-Refinement Diagnostic}

A4 studies a different executable path. Starting from the same grid support, it repeats deterministic local searches with radii
$\{0.45,0.28,0.18,0.11,0.07,0.045,0.03,0.02\}$ and recomputes LS at each depth up to $K=8$. If $q_k$ is NMSE in dB after depth $k$, the plateau gate stops at
\begin{equation}
  \widehat K=\min\{k\geq K_{\min}: q_{k-1}-q_k<\gamma\},
  \label{eq:plateau_gate}
\end{equation}
or at $K$ if the set is empty. This path has no trained CDL feature controller and should not be read as network depth. Its evidence supports only a high-load, 20 dB computation--quality trade-off under the tested schedule.
